% !PRINTED 154
\weberpage{154}
To compose three substitutions of the same dimension, $C$, $B$, $A$, one must introduce the expressions (10) for $z$ into a new linear substitution

\weq{12}{(u) = C(z)}

and express the variables $u$ in terms of the $x$:

\weq{13}{(u) = CBA(x).}

Since it is evidently immaterial whether one introduces the expressions (10) into (12), or whether one first expresses $z$ in terms of $y$ by (9) and then $y$ in terms of $x$ by (4), the \gesperrt{associative law} holds for this composition:

\weq{14}{C(BA) = (CB)A = CBA,}

which can also easily be confirmed by calculation, if one forms the elements of $C(BA)$ and $(CB)A$ according to (11). For both one finds the expression:

\[
\sum^i \sum^h c_i^{(k)}\, b_h^{(i)}\, a_l^{(h)}.
\]

We therefore state the theorem:

\par
3.\enspace\gesperrt{In the composition of the linear substitutions the associative law holds, but not always the commutative law.}

A multiplication in which all the multipliers are equal to one another, that is

\[
N = \begin{pmatrix} \nu, & 0, & \ldots & 0 \\ 0, & \nu, & \ldots & 0 \\ \cdot & \cdot & \cdot & \cdot \\ 0, & 0, & \ldots & \nu \end{pmatrix}
\]

shall be called a \gesperrt{similarity substitution}, and two substitutions which, like $A$ and $AN$ or $A$ and $NA$, can be derived from one another by composition with a similarity substitution, are called \gesperrt{similar substitutions}. From the law of composition the following theorems result at once:

\par
4.\enspace\gesperrt{A similarity substitution is permutable with every substitution $A$ of the same dimension; two similarity substitutions, compounded with one another, again give a similarity substitution, and two substitutions similar to a third are also similar to one another.}
